\section{convessità}

\subsection{insiemi convessi}

\begin{definition}[insieme convesso]
Sia $A \subset \RR^n$ un insieme. 
Diremo che $A$ è convesso se per ogni $x,y \in A$ 
si ha $[x,y]\subset A$ dove $[x,y]$ è il segmento che congiunge $x$ e $y$, cioè
\[
    [x,y] = \ENCLOSE{t x + (1-t) y \colon t \in [0,1]}.
\]

Diremo che $A$ è \emph{strettamente convesso} se per ogni $x,y \in A$, $x\neq y$ si ha
che $(x,y) \subset \stackrel\circ A$.
\end{definition}

\begin{theorem}[proiezione su un convesso]
Sia $A\subset \RR^n$ un insieme convesso e chiuso.
Dato $x\in \RR^n$ esiste un unico punto $y\in A$ tale che
\[
    \abs{x-y} = \min_{z\in A} \abs{x-z}.
\]
Possiamo quindi definire la proiezione su $A$, $P_A\colon \RR^n \to A$
che manda $x$ in $y$.

Allora $P_A$ è una funzione $1$-lipschiziana.
\end{theorem}
%
\begin{proof}
Per prima cosa osserviamo che fissato $y$ la funzione $x\mapsto \abs{x-y}$ 
è una funzione continua su $A$.
Inoltre se $A$ è illimitato tale funzione è anche coerciva. 
Quindi per il teorema di Weiestrass (generalizzato) sappiamo che esiste $y\in A$ 
che realizza il minimo.

Supponiamo ora che esista un altro punto $y'\in A$ che realizza il minimo.
Sia $r = \abs{x-y} = \abs{x-y'}$.
Visto che $\bar B_r(x)$ è strettamente convessa, 
il segmento $(y,y')$ è contenuto in $B_r(x)$.
E quindi, se fosse $y\neq y'$ 
esisterebbe un punto $z\in (y,y')$ che ha 
distanza da $x$ strettamente inferiore a $r$.
Ma per la convessità di $A$ l'intero segmento $[y,y']$ 
è contenuto in $A$ e quindi $z\in A$.
Questo è assurdo perché $r$ era la distanza minima
da $x$ a $A$.
Quindi c'è un unico punto $y\in A$ che realizza il minimo.

Per dimostrare che $P_A$ è $1$-lipschiziana, 
fissati $x, y \in \RR^n$ dobbiamo 
dimostrare che 
\begin{equation}\label{eq:convex4855}
  \abs{P_C(x) - P_C(y)} \le \abs{x-y}.
\end{equation}
Siano $p=P_C(x)$ e $q=P_C(y)$. 
Andremo ora a dimostrare che 
\begin{equation}\label{eq:convex4856}
  (x-p)\cdot (q-p) \le 0.
\end{equation}
Consideriamo il punto $p_t = (1-t) p + t q$.
Per la convessità di $C$ si ha $p_t \in C$ 
per ogni $t\in [0,1]$ e quindi 
la funzione $f(t) = \abs{x-p_t}^2$
ha minimo per $t=0$. 
Dunque deve essere $f'(0)\ge 0$. 
Ma  
\[ 
 f'(t) = 2 (x-p_t)\cdot q_t'
  = 2 (x-p_t)\cdot (q-p)
\]
e quindi $0\le f'(0) = 2 (x-p)\cdot (q-p)$
e si ottiene \eqref{eq:convex4856}.

Per concludere la dimostrazione osserviamo che
per simmetria si ha anche $(y-q)\cdot (p-q) \le 0$
e, sommando le due disuguaglianze, si ottiene 
dunque 
\[
 0 \ge (x - p - y + q)\cdot (q-p)
  = (x-y)\cdot (q-p) + \abs{q-p}^2
\]
da cui
\[
    \abs{q-p}^2 \le (y-x)\cdot (q-p) \le \abs{y-x}\cdot \abs{q-p}
\]
e quindi, dividendo per $\abs{q-p}$, si ottiene \eqref{eq:convex4855}.
\end{proof}

\begin{theorem}[piano di supporto]
  Sia $C\subset \RR^n$ un insieme chiuso e convesso e sia $p\in \partial C$.
  Allora esiste un semispazio $H$ con $H\supset C$ e $p\in \partial H$.
\end{theorem}
\begin{proof}
Visto che $p\in \partial C$ esiste $p_k\in \RR^n\setminus C$ con $p_k \to p$.
Grazie al teorema precedente possiamo considerare $q_k\in C$
il punto di minima distanza da $p_k$.
Se $r_k = \abs{p_k - q_k}$ allora 
$B_{r_k}(p_k) \cap C = \emptyset$ 
altrimenti ci sarebbe un punti di $C$ che dista da $p_k$ meno di $q_k$.
Consideriamo il semispazio chiuso $H_k$ con $\partial H_k$
tangente a $B_{r_k}(p_k)$ in $q_k$ e con
e $H_k \cap B_{r_k}(p_k) = \emptyset$.
Vogliamo dimostrare che $H_k \supset C$.

Supponismo per assurdo che esista $z\in C \setminus H_k$.
Allora il segmento $[q_k,z]$ da un lato è contenuto in $C$ 
per la convessità di $C$, d'altro lato
tale segmento dovrebbe entrare in $B_{r_k}(p_k)$  
(perché il complementare di $H_k$ è il tangente a $B_{r_k}(p_k)$ in $q_k$).
Ma questo è assurdo perché allora sul segmento $[q_k,z]$ ci sarebbe un punto di $C$ che dista da $p_k$ meno di $q_k$.
Dunque $H_k \supset C$.

Se consideriamo $v_k$ il versore normale a $\partial H_k$ in modo 
che $H_k = \ENCLOSE{x\in \RR^n\colon v_k\cdot (x-q_k) \le 0}$
allora visto che $v_k$ è limitata, 
per il teorema di Bolzano-Weierstrass 
possiamo estrarre una sottosuccessione $v_{k_j}$ che converge a
un certo $v$. 
Per conntinuità si ha $\abs{v_{k_j}} \abs v$ 
e quindi anche $v$ è un versore.
Visto che $H_k \supset C$, per ogni $x\in C$ si ha 
\[
  v_k \cdot (x-q_k)\le 0.
\]
Osservando che $\abs{q_k - p_k} \le \abs{p-p_k} \to 0$,
passando al limite nella disuguaglianza precedente otteniamo
\[
  v \cdot (x-p) \le 0
\]
che significa che $x\in H$ e $p\in \partial H$.
Per l'arbitrarietà di $x\in C$ otteniamo $C\subset H$, 
come volevamo dimostrare.
\end{proof}

\begin{corollary}
Un insieme chiuso $C\subset \RR^n$ è convesso 
se e solo se $C$ è intersezione (eventualmente infinita) 
di semispazi chiusi.
\end{corollary}
\begin{proof}
Da un lato è immediato verificare che l'intersezione di convessi 
è convessa (direttamente dalla definizione).
Viceversa $C$ è l'intersezione di tutti i semispazi di 
supporto, in quanto ogni $x\in \RR^n\setminus C$ è separato da $C$ 
dal semispazio di supporto passante dalla sua proiezione su $C$.
\end{proof}


\begin{definition}[funzione convessa]
Sia $A\subset \RR^n$ un insieme convesso e sia $f\colon A \to \RR$ una funzione definita su $A$.
Diremo che $f$ è convessa se per ogni $x,y \in A$ e per ogni $t \in [0,1]$ si ha
\begin{equation}\label{eq:convex}
    f(t x + (1-t) y) \leq t f(x) + (1-t) f(y).
\end{equation}

Diremo che $f$ è concava se $-f$ è convessa ovvero se vale la disuguaglianza 
opposta di \eqref{eq:convex}.
Diremo che $f$ è strettamente convessa se è convessa 
e \eqref{eq:convex} è verificata con disuguaglianza stretta per ogni $x,y \in A$, 
$x\neq y$ e per ogni $t \in (0,1)$.
Diremo che $f$ è strettamente concava se $-f$ è strettamente convessa.
\end{definition}

\begin{theorem}
Sia $A\subset \RR^n$ un insieme convesso e sia $f\colon A \to \RR$ una funzione.
Sia 
\[
  E_f \defeq \ENCLOSE{(x,y)\in \RR^{n+1}\colon y\ge f(x)}
\]
l'epigrafo di $f$.
Allora sono equivalenti:
\begin{enumerate}
    \item $A$ è convesso ed $f$ è convessa;
    \item $E_f$ è convesso.
\end{enumerate}
\end{theorem}

La convessità delle funzioni di più variabili può essere in buona parte 
ricondotta alla convessità delle funzioni di una variabile. 
Se $f\colon A\subset \RR^n \to \RR$, $x\in A$ e $v\in \RR^n$ 
possiamo definire la restrizione di $f$ alla retta passante per $x$ 
in direzione $v$, ponendo
\[
  g(t) = f(x+t v).
\]
Osserviamo quindi che $f$ è convessa se e solo se la corrispondente 
$g$ è convessa per ogni $x\in A$ e $v\in \RR^n$.

Si osserva inoltre che se $f$ è differenziabile si ha 
\[
  g'(t) = Df(x+t v)v
\]
e se $f$ è differenziabile due volte si ha
\[
  g''(t) = v^t D^2 f(x+t v) v.
\]

I tre seguenti risultati sono quindi conseguenza 
degli analoghi risultati per le funzioni di una variabile.

\begin{theorem}[criterio di convessità al primo ordine]
    Sia $A\subset \RR^n$ un insieme aperto convesso e sia $f\colon A \to \RR$.
    Se $f$ è convessa e $f$ è differenziabile in un punto $x_0 \in \stackrel\circ A$, 
    allora per ogni $x \in A$ si ha
    \begin{equation}\label{eq:conv474}    
        f(x) \ge f(x_0) + df(x_0)[x-x_0].
    \end{equation}
    Viceversa, se $f$ è differenziabile in $\stackrel\circ A$ 
    e \ref{eq:conv474} è verificata per ogni $x \in A$, allora $f$ è convessa.
\end{theorem}
\begin{theorem}[criterio di convessità al secondo ordine]
    Sia $A\subset \RR^n$ un insieme aperto convesso 
    e sia $f\colon A \to \RR$.
    Se $f$ è convessa e differenziabile due volte in un punto $x\in A$ 
    allora la matrice hessiana $D^2 f(x)$ è semidefinita positiva.
    Viceversa se $f$ è differenziabile due volte in $A$ 
    e per ogni $x\in A$ la matrice hessiana $D^2 f(x)$ è semidefinita positiva, 
    allora $f$ è convessa.
\end{theorem}


\begin{theorem}[combinazioni convesse]
    Sia $A\subset \RR^n$ un insieme convesso 
    e sia $f\colon A \to \RR$ una funzione convessa.
    Dati $x_1,\dots, x_N \in A$ e $t_1,\dots, t_N \in [0,1]$ con $\sum_{j=1}^N t_j = 1$ si ha
    \[
      f\enclose{\sum_{k=1}^N t_k x_k} \le \sum_{k=1}^N t_k f(x_k).
    \]
\end{theorem}
\begin{proof}
    Si dimostra per induzione su $N$.
    Se $N=1$ il risultato è banale. 
    Supponiamo dunque $N>1$ e supponiamo per induzione che il teorema sia vero
    per $N-1$ punti.
    Sia $s = 1-t_N$.
    Allora $s > 0$ e possiamo definire $s_k = t_k/s$ per $k=1,\dots, N-1$.
    Per ipotesi di induzione si ha 
    \[
      f\enclose{\sum_{k=1}^{N-1} s_k x_k} \le \sum_{k=1}^{N-1} s_k f(x_k).
    \]
    D'altra parte, per la convessità di $f$, si ha 
    \[
      f\enclose{\sum_{k=1}^N t_k x_k} = f\enclose{s \sum_{k=1}^{N-1} s_k x_k + t_N x_N}
      \le s f\enclose{\sum_{k=1}^{N-1} s_k x_k} + t_N f(x_N)
      \le s \sum_{k=1}^{N-1} s_k f(x_k) + t_N f(x_N)
      = \sum_{k=1}^N t_k f(x_k).
    \]
\end{proof}

\begin{theorem}[sottolivelli di una funzione convessa]
    Sia $A\subset \RR^n$ un insieme convesso e sia $f\colon A \to \RR$ una funzione convessa.
    Allora per ogni $c\in \RR$ gli insiemi \emph{sottolivello}
    \[
        \ENCLOSE{x\in A\colon f(x)\le c},
        \qquad \text{e} \qquad
        \ENCLOSE{x\in A\colon f(x) < c}
    \]
    sono convessi.
    Se inoltre $f$ è strettamente convessa, allora i sottolivelli sono strettamente convessi.
\end{theorem}
\begin{proof}
Segue direttamente dalla definizione di convessità.
\end{proof}

\begin{example}
Fissato $y\in \RR^n$
la funzione $f(x) = \abs{x-y}^2$ 
è strettamente convessa in quando $D^2 f(x) = 2 I$ è definita positiva. 
Dunque le palle $B_\rho(y)$, e $\bar B_\rho(y)$ essendo sottolivelli 
di tale funzione, sono strettamente convesse.
\end{example}

\begin{theorem}
Sia $A\subset \RR^n$ un insieme aperto convesso e sia $f\colon A \to \RR$ una funzione
strettamente convessa. Allora $f$ ha al più un punto di minimo globale.
\end{theorem}

\begin{theorem}[locale lipschitzianità delle funzioni convesse]
Sia $A\subset \RR^n$ un aperto convesso e sia $f\colon A \to \RR$ una funzione convessa.
Allora $f$ è localmente lipschitziana, cioè per ogni $x\in A$ esiste $\rho >0$ 
per cui $f$ è lipschitziana su $B_\rho(x)$.
In particolare $f$ è continua.
\end{theorem}
\begin{proof}
\emph{Passo 1:} mostriamo che $f$ è localmente superiormente limitata.
Dato $x\in A$ possiamo trovare un cubo $Q$ 
con centro in $x$ interamente contenuto in $A$.
Siano $x_1, \dots, x_{2^n}$ i vertici di $Q$.
Allora ogni punto del cubo è una combinazione convessa 
dei vertice e quindi il valore della funzione 
è inferiore alla combinazione convessa dei valori della funzione nei vertici
e quindi, in particolare, è inferiore al valore massimo 
sui vertici.

\emph{Passo 2:} mostriamo che $f$ è localmente inferiormente limitata.
Fissiamo un punto $x_0\in A$ e un cubo $Q$ con centro in $x_0$ 
su cui $f$ è superiormente limitata da una costante $M$.
Se $x\in Q$ allora anche $x' = 2 x_0 - x$ è in $Q$
e si ha $x_0 = (x+x')/2$. 
Quindi, per la convessità di $f$, si ha 
\[
  f(x_0) \le \frac 1 2 f(x) + \frac 1 2 f(x')
\]
da cui
\[
 f(x) \ge 2 f(x_0) - f(x') \ge 2f(x_0) - M.
\]

\emph{Conlusione.}
Dato $x_0\in A$ scegliamo $r>0$ 
per cui $m \le f(x) \le M$ per ogni $x\in \bar B_{2r}(x_0)$.
Dati $x,y \in B_r(x_0)$ se $y\neq x$ 
la semiretta uscente da $y$ passante per $x$ 
interseca la sfera $\partial B_{2r}(x_0)$ in un punto $z$.
Si avrà allora $x = (1-t) y + t z$ per un certo $t\in (0,1)$
e per la convessità di $f$ avremo 
\[
  f(x) - f(y) 
    \le (1-t) f(y) + t f(z) - f(y) 
    = t (f(z) - f(y))
    \le t (M-m). 
\]
Per stimare $t$, sappiamo che $\abs{x-y}\le 2r$ mentre $\abs{z-y} \ge r$ e quindi
\[
  t = \frac{\abs{x-y}}{\abs{z-y}} \le 2.
\] 
Dunque $f(x) - f(y) \le 2 (M-m) \abs{x-y}$.
Scambiando $x$ e $y$ otteniamo anche $f(y) - f(x) \le 2 (M-m) \abs{x-y}$
e quindi la conclusione.
\end{proof}
  
\begin{definition}[sottodifferenziale]
Sia $A\subset \RR^n$ un insieme aperto convesso e sia $f\colon A \to \RR$ una funzione convessa.
Dato $x \in A$ definiamo il \emph{sottodifferenziale} di $f$ in $x$ come l'insieme
\[
  \partial f(x) = \ENCLOSE{v\in \RR^n\colon f(y) \ge f(x) + v\cdot (y-x) \text{ per ogni } y\in A}.
\]
\end{definition}

La condizione al primo ordine per la convessità di $f$ 
garntisce che se $f$ è convessa e differenziabile in $x$ allora $\nabla f (x) \in \partial f(x)$.

\begin{theorem}[proprietà del sottodifferenziale]
Sia $A\subset \RR^n$ un insieme aperto convesso e sia $f\colon A \to \RR$ una funzione convessa.
Allora per ogni $x\in A$:
\begin{enumerate}
  \item $\partial f(x)$ è un insieme convesso;
  \item $\partial f$ è \emph{crescente} nel senso che
  per ogni $x,y\in A$, $v\in \partial f(x)$ e $w\in \partial f(y)$ si ha
  \[
    (v-w)\cdot (y-x) \ge 0;
  \]
  \item $\partial f$ è \emph{continuo} nel senso che se $x_k\to x$ e $v_k\in \partial f(x_k)$ con $v_k\to v$ allora $v\in \partial f(x)$.
  \item $f$ è differenziabile in $x$ se e solo se $\partial f(x)= \ENCLOSE{\nabla f(x)}$;
  \item $\partial f(x) \neq \emptyset$ per ogni $x\in A$.
\end{enumerate}
\end{theorem}
\begin{proof}
  Per mostrare che $\partial f(x)$ è convesso, fissiamo $v,w \in \partial f(x)$ e $t\in [0,1]$.
  Allora per ogni $y\in A$ si ha 
  \begin{align*}
    f(y) \ge f(x) + v\cdot (y-x) \\
    f(y) \ge f(x) + w\cdot (y-x)
  \end{align*}
  da cui
  \[
    f(y) = (1-t) f(y) + t f(y) 
      \ge f(x) + ((1-t)v + t w)\cdot (y-x)
  \]
  ovvero $(1-t)v + t w \in \partial f(x)$, 
  che è quanto dovevamo dimostrare.

  Per mostrare che $\partial f$ è crescente, fissiamo $v\in \partial f(x)$ e $w\in \partial f(y)$.
  Allora si ha 
  \begin{align*}
    f(y) \ge f(x) + v\cdot (y-x) \\
    f(x) \ge f(y) + w\cdot (x-y).
  \end{align*}
  Sommando le due disuguaglianze si ottiene il risultato.

  Per la continuità basta osservare che, grazie al teorema 
  precedente, $f$ è continua e quindi si può passare al limite 
  nella disuguaglianza che definisce $v_k\in \partial f(x_k)$.

  Per quanto riguarda la differenziabilità, abbiamo già osservato 
  che $\nabla f(x) \in \partial f(x)$ se $f$ è differenziabile in $x$.
  Ci rimane da dimostrare che se $\xi,\eta\in \partial f(x)$ 
  e $\xi\neq\eta$ allora $f$ non è differenziabile in $x$.
  Se $\xi\neq\eta$ allora esiste $v\in \RR^n$ 
  tale che $(\xi-\eta)\cdot v \neq 0$.
  Basta considerare la restrizione di $f$ alla retta passante per $x$ in direzione $v$.
  Sia $g(t) = f(x+t v)$, allora $g$ è convessa e
  \begin{align*}
   g(t) &= f(x+tv) \ge f(x) + \xi\cdot (tv) = g(0) + t \xi\cdot v\\
   g(t) &= f(x+tv) \ge f(x) + \eta\cdot (tv) = g(0) + t \eta\cdot v
  \end{align*}
  da cui $g(t) - g(0) \ge \max\ENCLOSE{t \xi\cdot v, t \eta\cdot v}$.
  Significa che $g$ non è derivabile in $0$ e quindi $f$ non è differenziabile in $x$.

  Per dimostrare che $\partial f(x) \neq \emptyset$ per ogni $x\in A$
  consideriamo il piano di supporto all'epigrafo di $f$ in $(x,f(x))$.
  Tale piano non può essere verticale, altrimenti 
  la funzione $f$ non sarebbe definita in un intorno di $x$.
  \end{proof}


  \begin{corollary}
  Sia $A\subset \RR^n$ un insieme aperto convesso e sia $f\colon A \to \RR$ una funzione convessa.
  Allora $f$ è estremo superiore di una famiglia di funzioni 
  lineari: 
  \[
   f(x) = \sup \ENCLOSE{f(z) + v\cdot (x-z) \colon z\in A, v\in \partial f(z)}.
  \]
  \end{corollary}

  \begin{example}[sottodifferenziale della norma]
  Sia $f(x) = \abs{x}$, $f\colon \RR^n \to \RR$. 
  Allora $\partial f(x) = \ENCLOSE{\frac{x}{\abs{x}}}$ se $x\neq 0$ e
  \[
    \partial f(0) = \bar B_1(0).
  \]
  Infatti se $x\neq 0$ allora $f$ è differenziabile in $x$ e $\nabla f(x) = \frac{x}{\abs{x}}$.
  Se $x=0$ allora $v\in \partial f(0)$ se e solo se $f(y) \ge v\cdot y$ per ogni $y\in \RR^n$.
  Ma $\abs{y} \ge v\cdot y$ per ogni $y\in \RR^n$ se e solo se $\abs{v}\le 1$.
  \end{example}
  \begin{example}[stadio]
  Sia $f\colon \RR^2\to \RR$ definita da 
  \[
    f(x,y) = \begin{cases}
     \sqrt{x^2+y^2} & \text{se } x\le 0, \\
     \abs{y} & \text{se } $x\in [0,1]$, \\
      \sqrt{(x-1)^2+y^2} & \text{se } x\ge 1.
    \end{cases}
  \]
  Allora se $x\in [0,1]$ si ha 
  \[
    \partial f (x,0) = \ENCLOSE{0} \times [-1,1].
  \]
  \end{example}